\section{Sistemi Lineari e Matrici}

\subsection{Sistemi Lineari}

\begin{exercisebox}[title={Esercizio 1: Parametri}]
Discutere, al variare del parametro $k \in \mathbb{R}$, le soluzioni del seguente sistema lineare:
$$
\begin{cases}
x + y + z = 1 \\
x + k y + z = 1 \\
x + y + k^2 z = k
\end{cases}
$$
\end{exercisebox}

\begin{solbox}
Scriviamo la matrice completa associata al sistema:
$$
(A|b) = \left(\begin{array}{ccc|c}
1 & 1 & 1 & 1 \\
1 & k & 1 & 1 \\
1 & 1 & k^2 & k
\end{array}\right)
$$
Applichiamo l'eliminazione di Gauss.
1. $R_2 \to R_2 - R_1$:
$$ \left(\begin{array}{ccc|c} 1 & 1 & 1 & 1 \\ 0 & k-1 & 0 & 0 \\ 1 & 1 & k^2 & k \end{array}\right) $$
2. $R_3 \to R_3 - R_1$:
$$ \left(\begin{array}{ccc|c} 1 & 1 & 1 & 1 \\ 0 & k-1 & 0 & 0 \\ 0 & 0 & k^2-1 & k-1 \end{array}\right) $$
La matrice è a scalini.
Discussione sui pivot ($k-1$ e $k^2-1$):
\begin{itemize}
    \item Se $k \neq 1$ e $k \neq -1$: Tutti i pivot sono non nulli. Soluzione unica.
    \item Se $k = 1$: Matrice $(0 \dots 0)$. Infinite soluzioni.
    \item Se $k = -1$: Terza riga diventa $0 \ 0 \ 0 \ | -2$. Impossibile.
\end{itemize}
\end{solbox}

\begin{exercisebox}[title={Esercizio 2: Sistema $2 \times 2$}]
Risolvere il sistema:
$$ \begin{cases} 2x - 3y = 7 \\ 4x + y = 21 \end{cases} $$
\end{exercisebox}

\begin{solbox}
Gauss:
$$ \left(\begin{array}{cc|c} 2 & -3 & 7 \\ 4 & 1 & 21 \end{array}\right) \xrightarrow{R_2 - 2R_1} 
\left(\begin{array}{cc|c} 2 & -3 & 7 \\ 0 & 7 & 7 \end{array}\right) $$
$7y = 7 \implies y = 1$.
$2x - 3(1) = 7 \implies 2x = 10 \implies x = 5$.
Soluzione: $(5, 1)$.
\end{solbox}

\begin{exercisebox}[title={Esercizio 3: Sistema Omogeneo}]
Determinare lo spazio delle soluzioni del sistema omogeneo:
$$ \begin{cases} x + 2y - z = 0 \\ 2x - y + 3z = 0 \end{cases} $$
\end{exercisebox}

\begin{solbox}
Matrice:
$$ \left(\begin{array}{ccc} 1 & 2 & -1 \\ 2 & -1 & 3 \end{array}\right) \xrightarrow{R_2 - 2R_1} 
\left(\begin{array}{ccc} 1 & 2 & -1 \\ 0 & -5 & 5 \end{array}\right) $$
La seconda riga corrisponde a $-5y + 5z = 0 \implies y = z$.
Sostituiamo nella prima: $x + 2(z) - z = 0 \implies x + z = 0 \implies x = -z$.
Ponendo $z = t$ (parametro libero), le soluzioni sono:
$$ \begin{pmatrix} x \\ y \\ z \end{pmatrix} = \begin{pmatrix} -t \\ t \\ t \end{pmatrix} = t \begin{pmatrix} -1 \\ 1 \\ 1 \end{pmatrix}, \quad \forall t \in \mathbb{R} $$
Lo spazio delle soluzioni ha dimensione 1 ed è generato da $v = (-1, 1, 1)$.
\end{solbox}
